\documentclass[11pt,a4paper]{ctexart}

% === Fonts and encoding ===
\setCJKmainfont{STHeiti}[BoldFont=STHeiti]
\setCJKsansfont{STHeiti}
\setCJKmonofont{STHeiti}

% === Page geometry ===
\usepackage[margin=1in]{geometry}

% === Tables ===
\usepackage{booktabs}
\usepackage{array}

% === Hyperlinks ===
\usepackage[hidelinks]{hyperref}
\urlstyle{same}

% === Captions ===
\usepackage{caption}
\captionsetup{font=small,labelfont=bf,skip=6pt}

% === Headers ===
\usepackage{fancyhdr}
\pagestyle{fancy}
\fancyhf{}
\fancyhead[L]{\small MSCE：多源一致性引擎}
\fancyhead[R]{\small\thepage}
\renewcommand{\headrulewidth}{0.3pt}

\title{\textbf{MSCE：多源一致性引擎\\——面向AI主张的对抗性交叉验证}}

\author{邓新航}
\date{2026年6月}

\begin{document}

\maketitle

\begin{center}
\small
\textbf{DOI:} \url{https://doi.org/10.5281/zenodo.20485437} \quad
\textbf{代码:} \url{https://github.com/sampson0826/msce} \quad
\textbf{邮箱:} \texttt{sampson1735937149@gmail.com}
\end{center}

\vspace{-2mm}

\begin{abstract}
单个AI模型会自信地给出错误答案，而一个模型无法可靠地检测自身错误——产生错误的内部表征，和用于自检的内部表征，是同一套。MSCE（多源一致性引擎）通过让六个前沿模型以刻意不同的推理策略回答同一问题，将输出经过三层对抗管道过滤，再由独立法官对照所有适用约束评估幸存答案，来解决上述问题。在总计386道题目的三个基准套件中，MSCE在前沿基准上答对\textbf{60/60}（GPT-5.5：51/60），在扩展基准上答对\textbf{114/120}（GPT-5.5：106/120），在综合基准上答对\textbf{180/206}（GPT-5.5：154/206）。全386题，MSCE平均准确率91.7\%对GPT-5.5的80.6\%，同时零个答案属于既错误又被赋予高置信度。我们将MSCE应用于两项真实验证任务：审查六个哈勃张力解决方案对照八个观测约束，以及交叉验证LeCun等人的LeWorldModel论文中的主张。MSCE不生成新内容。它是一个检查主张能否在同步、多角度审视下站住脚的工具——这件事单个模型和单个人类审稿人都不擅长。
\end{abstract}

\section{引言}

单个AI模型回答问题时，用户得到一个答案，通常还有一个置信度分数。模型无法标记自己的答案可能错误，因为生成答案的内部状态和生成置信度的是同一个。这并非罕见边角案例。在206道题目的基准上，GPT-5.5给出了\textbf{40个既错误又自评置信度高于0.7的答案}。其中包含被包装成确定结论的算术错误、被当作事实陈述的物理不可能论断、以及以流利但错误的论证包装的逻辑谬误。

当任务升级到人类也觉得困难的事情上，比如评审一篇科学论文，问题更严重。一篇论文提出五个断言。审稿人抓到两个问题。没有人——审稿人、作者、编辑——会同步检查每个断言对照所有相关约束：理论边界、观测数据、方法论假设、以及竞争团队的独立结果。那些在全量交叉检查下会失败的断言，之所以能通过同行评审，是因为交叉检查的认知负荷超出了人类工作记忆。

MSCE通过分离生成与验证来处理这个问题。六个模型独立回答同一问题，每个被强制走不同的推理路径。通过对抗过滤存活的答案，再对照明确声明的约束进行检验。一个独立的法官模型评估剩余的部分。核心设想很简单：六个以不同方式推理的模型就同一个错误答案达成一致的概率，远低于任何单一模型给出正确答案的概率。当它们确实一致时，那个一致性本身就是证据——不是证明，是证据——独立于任何一个模型的可靠性。

\section{相关工作}

集成方法如bagging和boosting通过聚合多个模型来提高预测效果。弱点是，在重叠数据上训练、架构相似的模型，往往共享盲点。对关联误差取平均并不能消除它们。

多智能体辩论让模型迭代地相互批评。这些系统可以收敛到口才而非真相——一个自信、表达力强的错误模型可以动摇不那么自信的正确模型。辩论的结果反映修辞技巧不亚于准确度。

LLM-as-judge让一个模型评价另一个模型。当法官和受评模型共享训练谱系时，法官继承了该模型家族的盲点。

共形预测对预测集提供分布性保证，但无法告诉你某个特定主张是否违反某个特定约束。

以上都不是验证架构。MSCE不是集成（模型不投票），不是辩论（模型在回答前不交互），不是单法官系统（判断需要约束满足，而非仅凭一个模型的意见）。

\section{架构}

\subsection{设计}

六个模型。六种推理策略。每个模型收到同一个问题，但被指示通过不同的认知路径推理：

\begin{table}[ht]
\centering
\caption{模型配置与异构认知策略}
\label{tab:models}
\begin{tabular}{>{\sffamily}ll>{\raggedright\arraybackslash}p{5.5cm}}
\toprule
\textbf{模型} & \textbf{策略} & \textbf{做什么} \\
\midrule
GPT-5.5 & 深度优先 & 沿一条推理链走到结论，再考虑替代方案 \\
Gemini 3.1 Pro & 广度优先 & 先列举合理路径，逐一评估，再选择 \\
Grok 4.1 & 反事实 & 先假设初始答案是错的，再逆向推理找到正确 \\
Kimi K2.5 & 直接推理 & 透明地逐步推理，每个中间结果显式呈现 \\
GPT-5.1 (thinking) & 科学深度 & 形成假设、推导可检验推论、评估证据、得出结论 \\
o4-mini & 约束传播 & 编码所有给定约束，传播逻辑蕴含，验证一致性 \\
\bottomrule
\end{tabular}
\end{table}

法官为Grok 4.1 (thinking)，选择它因为它来自与所有生成器不同的模型家族，减少自验证偏差。

\subsection{管道}

六个模型独立生成答案。答案经过三层过滤：

\begin{description}
\item[L1——置信度过滤。] 自报置信度低于0.3的答案被丢弃。模型自己都不确定，MSCE不覆盖它。

\item[L2——离群点检测。] 偏离答案簇语义中心超过$2\sigma$的答案被标记。一个既自信又孤立于群体的模型通常错了。注意一个失败模式：如果多数犯错，L2会丢弃少数正确答案。这是已知的取舍——MSCE优先检测自信错误而非保留少数正确，这对验证适用但对发现不适用。

\item[L3——盲点检查。] 当所有模型高置信度一致时，MSCE估算集体盲点风险。在相似数据和目标上训练的模型可能共享盲点。高一致不保证正确。L3不拒绝答案；当一致性高但约束覆盖低时，它挂一个警告。
\end{description}

过滤后，存活的答案进入交叉验证矩阵。每个（答案，约束）对被归类为通过、紧张或违反。法官评估矩阵，输出最终答案及校准置信度。

\subsection{交叉验证矩阵}

对每个主张，MSCE对照所有适用约束进行检验。输出是一个矩阵而非单一数字：

\begin{itemize}
    \item \textbf{通过：} 主张与约束一致。
    \item \textbf{紧张：} 主张触及约束边界但不打破。
    \item \textbf{违反：} 主张与约束矛盾。
\end{itemize}

当约束本身冲突时（如Planck $H_0 = 67.4$与SH0ES $H_0 = 73.0$），MSCE将约束对标记为内部不一致，并在两种配置下分别报告结果——关闭一个冲突约束、关闭另一个、两个都关闭。这避免了主张仅因约束相互矛盾就被标记为"违反所有约束"。

\subsection{置信度校准}

MSCE置信度不是模型的自我报告概率。它从三个信号计算：

\begin{enumerate}
    \item 幸存模型的一致率。
    \item 通过的约束与总约束之比。
    \item 模型间分歧的程度。
\end{enumerate}

结果系统性低于单个模型的自我置信度。GPT-5.5报告0.75的平均置信度而错误率15\%。MSCE报告0.49—0.57的平均置信度而错误率8.3\%。更低的置信度在这里是优点——意味着系统知道什么时候不确定。

我们将MSCE置信度超过0.8的答案归为"高置信度"。在此阈值下，MSCE在全386题上产生\textbf{零个高置信度错误}。这个阈值是保守的：MSCE的置信度很少超过0.8。在综合基准上MSCE错了26题；这些错误中最高置信度为0.62。

关于阈值比较：全文中GPT-5.5的"高置信度错误"使用自报置信度$>$0.7作为阈值。这种不对称（MSCE用0.8，GPT-5.5用0.7）是因为两个置信度量表校准不同。MSCE置信度是合成值，通常跨度0.3—0.7（系统平均：0.49—0.57）。GPT-5.5自报置信度平均0.74—0.75，且频繁超过0.9。将GPT-5.5的阈值提高到$>$0.8会减少但不会消除其高置信度错误数；模式（MSCE零，GPT-5.5非零）在任何合理阈值下都成立。

\section{基准}

\subsection{设计}

四个基准套件，在所有组成模型的最新训练截止日之后构建，避免数据污染：

\begin{table}[ht]
\centering
\caption{基准套件设计}
\begin{tabular}{lrl}
\toprule
\textbf{套件} & \textbf{题目数} & \textbf{领域} \\
\midrule
Pilot（仅调参） & 20 & 数学、逻辑、科学、语言 \\
Frontier & 60 & 数学(10)、逻辑(10)、科学(20)、语言(20) \\
Expanded & 120 & 以上 + 约束传播(9)、跨域(10) \\
Comprehensive & 206 & 全部6个领域，L1--L4难度梯度 \\
\bottomrule
\end{tabular}
\end{table}

20题的Pilot套件用于管道参数调优，排除在所有汇总指标之外。以下所有数据来自三个主套件（$60 + 120 + 206 = 386$题）。

\subsection{汇总结果}

全三个主套件（386题）：

\begin{table}[ht]
\centering
\caption{各领域准确率对比（三基准汇总）}
\begin{tabular}{lrrrr}
\toprule
\textbf{领域} & \textbf{题数} & \textbf{MSCE} & \textbf{GPT-5.5} & \textbf{差距} \\
\midrule
数学 & 60 & 98.3\% & 95.0\% & $+3.3$pp \\
科学 & 87 & 98.9\% & 73.6\% & $+25.3$pp \\
跨域 & 43 & 86.0\% & 60.5\% & $+25.5$pp \\
逻辑 & 57 & 94.7\% & 89.5\% & $+5.2$pp \\
约束传播 & 52 & 71.2\% & 61.5\% & $+9.7$pp \\
语言 & 87 & 93.1\% & 93.1\% & $0.0$pp \\
\midrule
\textbf{合计} & \textbf{386} & \textbf{91.7\%} & \textbf{80.6\%} & \textbf{$+11.1$pp} \\
\bottomrule
\end{tabular}
\end{table}

最大差距出现在科学和跨域推理——这两个领域从基于约束的验证中受益最多。语言成绩完全相同，这说得通：语言推理涉及的显式约束较少。

关于比较的说明：MSCE使用六个模型加一个法官。与其比较的是单次GPT-5.5调用，即$7\times$计算量对$1\times$。我们报告这个比较是因为单模型推断是大多数用户的默认部署方式。第6.3节讨论了计算量问题，以及需要哪些消融实验来将收益归因于架构而非规模。简言之：在没有消融实验的情况下，11.1pp的差距是MSCE架构贡献的上界，而非精确度量。

\subsection{Frontier基准（60题）}

MSCE答对60/60。Wilson二项置信区间（95\%）：[94.0\%, 100\%]。GPT-5.5答对51/60（85.0\%）。MSCE在此集上的平均置信度为0.575；GPT-5.5为0.747。GPT-5.5产生了6个自报置信度超0.7但错误的答案。

\subsection{Expanded基准（120题）}

MSCE：114/120（95.0\%）。GPT-5.5：106/120（88.3\%）。平均置信度：MSCE 0.511，GPT-5.5 0.743。GPT-5.5有8个高置信度（$>$0.7）错误。

\subsection{Comprehensive基准（206题）}

MSCE：180/206（87.4\%）。GPT-5.5：154/206（74.8\%）。平均置信度：MSCE 0.486，GPT-5.5 0.744。GPT-5.5有40个高置信度错误；MSCE没有超过0.62置信度的错误。此基准使用L1--L4难度梯度，旨在压测推理深度而非广度。

\subsection{对抗鲁棒性}

80题分四类陷阱：

\begin{table}[ht]
\centering
\caption{对抗鲁棒性测试}
\begin{tabular}{llr}
\toprule
\textbf{陷阱} & \textbf{测试什么} & \textbf{MSCE} \\
\midrule
逻辑矛盾 & 自驳前提 & 20/20 \\
虚假前提 & 嵌入的事实错误 & 20/20 \\
权威陷阱 & 被包装为共识的虚假命题 & 20/20 \\
提示注入 & 跳过推理的指令 & 19/20 \\
\midrule
\textbf{合计} & & \textbf{79/80（98.8\%）} \\
\bottomrule
\end{tabular}
\end{table}

唯一失败来自一道将注入嵌入代码块的题目——约束传播模型将其解析为数据而非指令，产生了法官错误解决的分歧。

\subsection{知识边界校准}

30题分三个可知性层级：

\begin{table}[ht]
\centering
\caption{知识边界校准}
\begin{tabular}{llrr}
\toprule
\textbf{层级} & \textbf{描述} & \textbf{MSCE} & \textbf{GPT-5.5} \\
\midrule
第1层 & 已知事实 & 10/10 & 10/10 \\
第2层 & 真正不确定 & 6/10 & 0/10 \\
第3层 & 当前不可知 & 4/10 & 0/10 \\
\bottomrule
\end{tabular}
\end{table}

GPT-5.5对第2、3层全部20题赋予了$>$0.7的置信度——并且全部答错。MSCE在大多数情况下正确识别了自己不知道，尽管它仍在6道第3层题目上尝试了不该尝试的回答。

\section{案例研究}

\subsection{哈勃张力：六个方案，八个约束}

局部（SH0ES: $H_0 \approx 73.0$）与早期宇宙（Planck: $H_0 \approx 67.4$）膨胀率测量之间的$5\sigma$差异催生了数百个解决方案。我们将六个主要方案类别通过MSCE对照八个独立观测约束：

\begin{table}[ht]
\centering
\caption{哈勃张力方案 vs 八个观测约束}
{\small
\begin{tabular}{lccccccccc}
\toprule
\textbf{方案} & \textbf{得分} & \textbf{CMB} & \textbf{BAO} & \textbf{SN} & \textbf{BBN} & $\mathbf{S_8}$ & \textbf{Age} & \textbf{Grav} & \textbf{X} \\
\midrule
衰减暗物质 & 0.358 & P & P & T & P & P & P & P & T \\
额外中微子 & 0.287 & P & T & P & V & T & P & P & T \\
修正引力 & 0.253 & P & P & T & P & T & T & V & T \\
局域空洞 & 0.171 & P & T & P & P & P & P & P & T \\
系统误差 & 0.108 & P & P & P & P & P & P & P & P \\
早期暗能量 & 0.076 & V & V & P & P & V & P & P & T \\
\bottomrule
\end{tabular}}
\par\smallskip
P = 通过，T = 紧张，V = 违反，X = 跨约束一致性。
\end{table}

没有任何单一方案满足全部八个约束。没有任何配对组合满足。

"系统误差"条目值得说明：它不是物理模型，不做出可证伪的预测。它之所以通过所有约束，是因为没有提出任何足够具体到能被违反的断言。评分反映了这一点——它得了第二低分（0.108），尽管零违反。EDE排名最末。它同时违反CMB、BAO和$S_8$——三个观测上相互独立的约束。

一个限制：约束集同时包含Planck和SH0ES，它们的$H_0$值本身处于紧张状态。我们通过分别关闭各个冲突约束、再两个都关闭来处理。模式（EDE最末、无方案通过所有约束）在所有配置下保持。

\subsection{LeCun LeWorldModel审计}

LeCun等人提出了一个从像素端到端训练的JEPA世界模型。论文提出了五个中心主张。我们将每个分解为可验证的命题并附加独立约束：论文自报的局限性、同期结果（Sub-JEPA，2026年5月）、以及主张的数学基础。

\begin{table}[ht]
\centering
\caption{MSCE审计LeCun LeWorldModel：五条主张，五条否定}
\begin{tabular}{>{\raggedright\arraybackslash}p{4.5cm}lrr}
\toprule
\textbf{主张} & \textbf{裁决} & \textbf{模型一致} & \textbf{置信度} \\
\midrule
SIGReg"可证明最优"防坍塌 & FALSE & 4/5 & 0.57 \\
48倍加速 vs DINO-WM & UNFAIR & 4/5 & 0.57 \\
仅需1个超参数 & FALSE & 4/5 & 0.57 \\
JEPA坍塌"已解决" & FALSE & 4/4 & \textbf{0.95} \\
潜表示涌现物理理解 & FALSE & 5/5 & 0.57 \\
\bottomrule
\end{tabular}
\end{table}

第四个主张的非同寻常的高置信度（0.95）来自约束集非同寻常的紧密：论文自报的测试环境（四个低维控制任务）、一项报告世界模型脆性的同期基准研究、以及缺乏与非PLDM防坍塌方法的对比。当约束如此具体且相互印证时，模型间一致性上升。

我们也审计了Google DeepMind的Co-Scientist（Nature，2026年5月），完整结果见代码仓库。

\section{讨论}

\subsection{为什么有效}

引擎有效是因为推理策略差异足够大，导致它们的失败模式不完全相关。一个锁定在错误前提上的深度优先推理者犯的错误，与一个遗漏约束的约束传播推理者犯的错误不同。当它们通过不同路径到达同一答案时，该答案经历了独立检验。

这也就是容错硬件中$N$模冗余的原理：独立单元很少以相同方式同时失效。这里的独立性是近似的——所有六个模型都基于Transformer、在重叠数据上训练——但策略提示将它们推向不同的推理方向。

\subsection{有意义的置信度}

AI系统最危险的输出不是错误答案，是自信地给出的错误答案。MSCE的置信度数值与单一模型置信度行为不同：平均更低，但高时更准。

因为MSCE置信度并非从token概率分布中提取。它从可观测信号计算：模型一致吗？约束通过了吗？收敛前分歧多大？这些信号外在于任何单一模型的内部状态，因此与实际正确率的相关性更好。

代价是覆盖范围。MSCE拒答（置信度$<$0.3）的频率高于单一模型。对于验证用例——论文审计、主张检查、安全测试——拒答优于虚假自信。

\subsection{局限性}

\textbf{不公平基线。} MSCE用六个模型加法官对照单一模型基线。11.1pp的差距应被解读为上界。需要消融实验——六个模型简单投票、GPT-5.5 self-consistency@7、单模型准确率与集成贡献对比——来分离MSCE的架构贡献与多算力的原始效应。

\textbf{认知异构性未度量。} 策略被描述为异构的，但这是设计意图，不是已验证的属性。

\textbf{可复现性。} 组成模型为闭源API，可能变更或下架。精确复现将不可能。

\textbf{L2可能丢弃正确的少数答案。} 如果多数共享盲点，少数正确答案会被移除。

\textbf{约束选择偏差。} MSCE使约束集显式可见，但约束选择仍然是人的决定。

\textbf{无人类基线。} 与人类审稿人的受控对比尚未进行。

\section{结论}

MSCE是一个验证工具。它让六个模型互相检验，过滤掉无法通过对抗审视的答案，并对剩余的部分对照显式约束做检查。跨越386题，它答对91.7\%——比最佳单一模型高11.1个百分点——同时零个答案是错且自信的。在大多数答错的情况下，它标记了自己的不确定。

MSCE就做一件事。它不生成、不创造、不发现。它检验。在AI生成的主张变得丰富且廉价的世界里，检验才是瓶颈。MSCE是一次拓宽它的尝试。

\end{document}
